% chapters/appendix_a_proofs.tex

This appendix collects the detailed derivations referenced throughout the main
text: the algebraic manipulations underlying the Gronwall closure in Theorem~3.1
(Chapter 3), the discretisation error lemma used in Theorem~4.1 (Chapter 4), and
the Riccati system derivations underlying the closed-form solutions of Chapters 7
and 8.


\section{Detailed Derivation of the Contraction Estimate (Chapter 3)}
\label{sec:appA-contraction}

This section provides the algebraic detail omitted from \Cref{sec:wp-stability} of
Chapter 3, specifically the derivation of the control-difference Lipschitz constant
$C_\alpha$, the BSDE estimate constant $C_{\mathrm{BSDE}}$, and the combination
step yielding the differential inequality \eqref{eq:gronwall-differential}.

\subsection{Derivation of $C_\alpha$}

Recall the optimality condition $\partial_\alpha H(t,X_t,\mu_t,p_t,\alpha_t^*) = 0$.
Differentiating implicitly with respect to $(x, p, \mu)$ and using the uniform
concavity Assumption~3.5 ($\partial_{\alpha\alpha} H \leq -\lambda_H < 0$), the
implicit function theorem gives
\[
    \nabla_{(x,p,\mu)} \alpha^* = -\big(\partial_{\alpha\alpha} H\big)^{-1}
    \nabla_{(x,p,\mu)} \partial_\alpha H.
\]
Since $\|\partial_{\alpha\alpha} H\|^{-1} \leq 1/\lambda_H$ and
$\|\nabla_{(x,p,\mu)} \partial_\alpha H\| \leq L_b + L_\mu$ (by the Lipschitz
properties of $b$ and the bound on the Lions derivative of $f$), we obtain
\[
    C_\alpha = \frac{L_b + L_\mu}{\lambda_H}.
\]
This is the constant appearing in \eqref{eq:control-lipschitz}.

\subsection{Derivation of $C_{\mathrm{BSDE}}$}

The linear BSDE \eqref{eq:adjoint-difference-bsde} for $\Delta p_t$ has driver
coefficients $A_t, B_t, D_t$ bounded by the second derivatives of $H$, which are in
turn bounded by Assumption~3.3 (regularity). Standard linear BSDE estimates (El
Karoui, Peng, and Quenez, 1997, Proposition 2.1) give
\[
    \mathbb{E}\Big[\sup_t |\Delta p_t|^2\Big] \leq C \,\mathbb{E}\Big[|\Delta X_T|^2
    + \int_0^T \big(|\Delta X_t|^2 + |\alpha_t^1-\alpha_t^2|^2 +
    W_2(\mu_t^1,\mu_t^2)^2\big)\, dt\Big],
\]
where $C$ depends on $T$, the bounds on $A_t, B_t, D_t$, and --- crucially --- on
$\|\sigma_0\|_F^2$ through the martingale representation term $\Delta r_t$. The
common-noise contribution enters because the quadratic variation of the
common-noise martingale term in the BSDE energy identity is
$\mathbb{E}[\int_0^T |\Delta r_t|^2\, dt]$, and the martingale representation
theorem bounds this by the same right-hand side scaled by $\|\sigma_0\|_F^2$. This
gives $C_{\mathrm{BSDE}} = C_0' (1 + \|\sigma_0\|_F^2)$ for a constant $C_0'$
depending only on $T$ and the Lipschitz constants.

\subsection{Combining the Estimates}

Substituting \eqref{eq:control-lipschitz} into \eqref{eq:gronwall-setup} and
\eqref{eq:bsde-estimate}, and writing $y_t = \mathbb{E}[|\Delta X_t|^2] +
\mathbb{E}[|\Delta p_t|^2]$, the two coupled inequalities
\begin{align*}
    \frac{d}{dt}\mathbb{E}[|\Delta X_t|^2] &\leq (2L_b + L_b^2)\,
        \mathbb{E}[|\Delta X_t|^2] + L_b\, \mathbb{E}[W_2(\mu_t^1,\mu_t^2)^2] +
        L_b C_\alpha^2 \big( \mathbb{E}[|\Delta X_t|^2] + \mathbb{E}[|\Delta p_t|^2]
        + \mathbb{E}[W_2(\mu_t^1,\mu_t^2)^2] \big), \\
    \mathbb{E}[|\Delta p_t|^2] &\leq C_{\mathrm{BSDE}}\, \mathbb{E}\Big[|\Delta
        X_T|^2 + \int_0^T \big( y_s + \mathbb{E}[W_2(\mu_s^1,\mu_s^2)^2]
        \big)\, ds\Big]
\end{align*}
can be combined (using Gronwall's inequality on the second to bound
$\mathbb{E}[|\Delta p_t|^2]$ pointwise in terms of an integral of $y_s$, then
substituting into the first) into the single differential inequality
\[
    \frac{d}{dt} y_t \leq C_1\, y_t + C_2\, \mathbb{E}[W_2(\mu_t^1,\mu_t^2)^2],
\]
with $C_1 = 2L_b^2 + 2\|\sigma_0\|_F^2 + C_{\mathrm{reg}}$ and
$C_2 = L_b + C_\alpha L_b + C_{\mathrm{BSDE}} L_\mu$, where $C_{\mathrm{reg}}$
collects the remaining Lipschitz/regularity terms from $C_\alpha$ and
$C_{\mathrm{BSDE}}$ not already accounted for in the leading $2L_b^2$ and
$2\|\sigma_0\|_F^2$ terms. This is precisely \eqref{eq:gronwall-differential}.


\section{Discretisation Error Lemma (Chapter 4)}
\label{sec:appA-discretisation}

This section proves the discretisation error bound used in the proof of
Theorem~\ref{thm:numerical-convergence} (Lemma~C.9 in the original
chapter-internal numbering).

\begin{lemma}[Discretisation error of the fixed-point map]
\label{lem:appA-discretisation-error}
Let $\Phi$ be the continuous-time MFG fixed-point map (Chapter 3) and
$\Phi^\Delta$ its discrete analogue (Chapter 4, \Cref{sec:na-discretisation}).
Under Assumptions~1.13--1.19 and the smoothness Assumption~4.1, for any measure
flow $\mu \in \mathcal{M}$,
\[
    \|\Phi^\Delta(\mu) - \Phi(\mu)\|_{\lambda,T} \leq C_{\mathrm{disc}}(\Delta
    t^{1/2} + \Delta x),
\]
where $C_{\mathrm{disc}}$ depends on the Lipschitz/smoothness constants but not on
$\mu$, $\Delta t$, or $\Delta x$.
\end{lemma}

\begin{proof}
Fix $\mu \in \mathcal{M}$. By definition, $\Phi(\mu)$ is obtained by solving the
continuous-time agent FBSDE \eqref{eq:agent-fbsde} exactly, while $\Phi^\Delta(\mu)$
is obtained from the discrete scheme
\eqref{eq:discrete-forward-sde}--\eqref{eq:discrete-fp-spde}. We bound the
difference in three steps, corresponding to the three discretised components.

\textit{Step 1 (forward SDE).} By Lemma~\ref{lem:weak-error-forward}, the weak
error between the continuous state $X_t$ and discrete state $X^n$ (at matching
times $t_n = n\Delta t$) is $O(\Delta t)$ for smooth test functions. Combined with
Lemma~\ref{lem:weak-error-propagation}, this gives
$\mathbb{E}[|X^n - X_{t_n}|^2] \leq C\Delta t$, hence a strong-error contribution
of order $\Delta t^{1/2}$ to the Wasserstein distance between laws.

\textit{Step 2 (adjoint BSDE).} By Lemma~\ref{lem:weak-error-adjoint}, the
discrete adjoint process satisfies $|\mathbb{E}[\psi(p^n)] -
\mathbb{E}[\psi(p_{t_n})]| \leq C\Delta t$ for smooth $\psi$. This propagates to
the optimal control via the Lipschitz dependence $\alpha^* = p/\lambda$
(\Cref{eq:alpha-optimal-control}), contributing a further $O(\Delta t)$ term, which
after the same square-root argument as Step 1 contributes $O(\Delta t^{1/2})$ to
the Wasserstein bound.

\textit{Step 3 (Fokker--Planck SPDE / measure consistency).} By
Theorem~\ref{thm:fd-convergence}, the finite-difference (or particle) discretisation
of the measure flow satisfies
$\sup_n \mathbb{E}[\|\rho^n - \rho_{t_n}\|_{L^2}^2]^{1/2} \leq C(\Delta t^{1/2} +
\Delta x)$. Since the $L^2$ density norm controls the Wasserstein-2 distance for
measures with bounded density ratios (a consequence of the
Cauchy--Schwarz/transport-cost duality; see \citep{villani2009}, Remark 6.6), this
contributes the dominant $O(\Delta t^{1/2} + \Delta x)$ term.

Summing the three contributions (each bounded uniformly over $\mu$ by the
Lipschitz continuity of the coefficients, which ensures the discretisation
constants do not depend on which measure flow $\mu$ is plugged into $\Phi$ versus
$\Phi^\Delta$) gives
\[
    \|\Phi^\Delta(\mu) - \Phi(\mu)\|_{\lambda,T} \leq C_{\mathrm{disc}}(\Delta
    t^{1/2} + \Delta x),
\]
as claimed, with $C_{\mathrm{disc}}$ depending only on the Lipschitz and
smoothness constants of $b, f, g$ and the CFL-condition constants, uniformly over
bounded sets of measure flows.
\end{proof}


\section{Derivation of the Riccati Systems (Chapters 7 and 8)}
\label{sec:appA-riccati}

This section provides the detailed derivation of the Riccati equations
\eqref{eq:sr-riccati-quadratic}--\eqref{eq:sr-riccati-linear} (Chapter 7) and
\eqref{eq:alpha-riccati-system} (Chapter 8), which share the same derivation
structure since both arise from the linear-quadratic MFG ansatz.

\subsection{Setup}

We work with the scalar version (Chapter 7's systemic risk model; the matrix
version of Chapter 8 follows by an identical argument with matrix-valued
coefficients and the same coefficient-matching strategy, replacing scalar products
by the appropriate matrix products). Substituting the quadratic ansatz
\eqref{eq:sr-value-ansatz} into the HJB equation
\[
    \partial_t u + \frac{\sigma^2}{2}\partial_{xx} u + \frac{\sigma_0^2}{2}
    \partial_{xx} u + H^*(t,x,\mu,\partial_x u) + \int_{\mathbb{R}}
    \frac{\delta u}{\delta\mu}(\mu)(y) \cdot b(t,y,\mu,\alpha^*)\, \mu(dy) = 0,
\]
and matching coefficients of $x^2$, $x\bar\mu$, $\bar\mu^2$, $x$, $\bar\mu$, and
constants yields the system of ODEs. The Lions derivative of $u$ with respect to
$\mu$ at $y$ is
\[
    \frac{\delta u}{\delta\mu}(\mu)(y) = -B_t x - C_t \bar\mu - E_t.
\]
Substituting into the integral term and using $\int y\, \mu(dy) = \bar\mu$ gives
the contributions to the coefficients. The terms involving $D_t, E_t, F_t$ arise
from the linear part of the Hamiltonian and the bailout policy $b_t$ (Chapter 7) or
the baseline drift (Chapter 8, where $d_t = e_t = f_t = 0$ since Example~1.1 has no
exogenous linear forcing term).

\subsection{Quadratic Terms}

\textbf{Correction (this revision):} the derivation below replaces an earlier
version that, despite an empirical optimality check confirming the
\emph{sign} of the $A_tB_t/\lambda$ cross term in $\dot{B}_t$, was found on
full re-derivation to (i) contain an independent sign error in $\dot{A}_t$
unrelated to that check, and (ii) omit a Lions-derivative (master-equation)
coupling term from $\dot{B}_t$, $\dot{D}_t$, $\dot{C}_t$, $\dot{E}_t$, and
$\dot{F}_t$ alike. Both issues are corrected together here, since they are
intertwined: the corrected $A_t$ feeds into every other coefficient. The
correction was found by (a) re-deriving the Hamiltonian directly from the
cost functional and adjoint BSDE of Section~\ref{sec:sr-stackelberg} rather
than from the previously stated $H^*$ alone, which pins down the maximised
Hamiltonian and its terminal condition unambiguously, and (b) deriving the
mean-field coupling term explicitly from the same conditional-mean
consistency argument used for the aggregate drift $\bar{b}(t,\bar\mu)$
below. The full corrected system was verified three independent ways:
(1) direct symbolic substitution showing the PDE residual vanishes
identically; (2) numerical integration with two independent solvers
(LSODA and Radau) agreeing to all displayed digits; (3) the corrected
$A_t,B_t,D_t$ equations, when stripped of the mean-field coupling term,
exactly reproduce the maximised-Hamiltonian-only matching below, confirming
internal consistency of the two-step derivation.

The cost functional (Eq.~\ref{eq:sr-running-cost} and surrounding text) and
adjoint BSDE (Eq.~\ref{eq:sr-adjoint-bsde}) give the Hamiltonian
\[
    H(t,x,\mu,p,\alpha) = p\big[\kappa(\bar\mu-x)+\alpha+b_t\big] -
    \frac{c_1}{2}x^2 - \frac{\lambda}{2}\alpha^2,
\]
maximised over $\alpha$ (the bank maximises negative cost) at $\alpha^*=p/\lambda$,
giving the maximised Hamiltonian
\[
    H^*(t,x,\mu,p) = p\cdot\kappa(\bar\mu-x) + \frac{p^2}{2\lambda} + pb_t -
    \frac{c_1}{2}x^2,
\]
identical in form to the previously stated $H^*$. With $p=\partial_x V=-A_tx-B_t\bar\mu-D_t$
(confirmed via the adjoint terminal condition $p_T=-c_2X_T$, which matches
$\partial_xV(T,x,\mu)=-c_2x$ exactly when $A_T=c_2,B_T=D_T=0$), matching
coefficients of $x^2$ in $\partial_tV+H^*=0$ gives
\[
    \dot{A}_t = 2\kappa A_t + \frac{A_t^2}{\lambda} - c_1, \qquad A_T = c_2.
\]
\textbf{This differs from the previously stated equation in the sign of both
the $A_t^2/\lambda$ and $c_1$ terms.} The corrected equation is the standard
sign for a cost-\emph{minimisation} Riccati equation (equivalently, a
value-\emph{maximisation} equation written with $p=\partial_xV$), and was
independently verified by formulating the bank's problem as a textbook
two-state ($x,\bar\mu$) linear-quadratic regulator and solving the resulting
matrix Riccati equation directly, which reproduces this same $A_t$ with no
free sign choices. A welcome consequence: under this corrected equation, the
numerical illustration of Section~\ref{sec:sr-numerical} no longer exhibits
the finite-time blow-up reported in an earlier revision, and the originally
intended parameter $c_2=5.0$ may be restored without modification (see the
updated Table~\ref{tab:sr-baseline-parameters} and discussion in
Section~\ref{sec:sr-numerical}).

The full master equation for $V(t,x,\mu)=-\tfrac12A_tx^2-B_tx\bar\mu-\tfrac12C_t\bar\mu^2
-D_tx-E_t\bar\mu-F_t$ also requires the Lions-derivative coupling term,
which (since $\delta V/\delta\mu(y)=-B_tx-C_t\bar\mu-E_t$ is independent of
$y$) reduces to $\bar{b}(t,\bar\mu)\cdot\partial_{\bar\mu}V$, where
$\bar{b}(t,\bar\mu)=\mathbb{E}_{y\sim\mu}[\kappa(\bar\mu-y)+\alpha^*(t,y,\mu)+b_t]$
is the equilibrium aggregate drift of $\bar\mu_t$ itself. Since
$\alpha^*(t,y,\mu)=p(t,y,\mu)/\lambda$ is linear in $y$, this expectation
needs only the first moment $\mathbb{E}[y]=\bar\mu$ (no higher-moment terms
enter), giving exactly
\[
    \bar{b}(t,\bar\mu) = -\frac{(A_t+B_t)\bar\mu + D_t}{\lambda} + b_t.
\]
Together with the common-noise diffusion terms
$\tfrac12(\sigma^2+\sigma_0^2)\partial^2_{xx}V+\sigma_0^2\partial^2_{x\bar\mu}V
+\tfrac12\sigma_0^2\partial^2_{\bar\mu\bar\mu}V$ -- present because $\bar\mu_t$
is itself a stochastic process driven by the same common noise $\sigma_0$ as
$X_t$ -- the full master-equation HJB is
\[
    \partial_tV + H^*(t,x,\mu,\partial_xV) + \bar{b}(t,\bar\mu)\partial_{\bar\mu}V
    + \tfrac12(\sigma^2+\sigma_0^2)\partial^2_{xx}V + \sigma_0^2\partial^2_{x\bar\mu}V
    + \tfrac12\sigma_0^2\partial^2_{\bar\mu\bar\mu}V = 0.
\]
The coupling term $\bar{b}\cdot\partial_{\bar\mu}V$ contributes \emph{only}
to the $x\bar\mu$, $\bar\mu^2$, $x$, $\bar\mu$, and constant coefficients
(verified symbolically: its contribution to the $x^2$ coefficient is
identically zero, so $A_t$ above is unaffected by this correction).
Matching coefficients of $x\bar\mu$ now gives
\[
    \dot{B}_t = \kappa B_t - \kappa A_t + \frac{2A_tB_t}{\lambda} +
    \frac{B_t^2}{\lambda}, \qquad B_T = 0,
\]
which agrees with the previously corrected sign of the $A_tB_t/\lambda$ term
(confirming that empirical check remains valid) but adds the
$2A_tB_t/\lambda+B_t^2/\lambda$ master-equation contribution, replacing the
single $A_tB_t/\lambda$ term that appeared before. Matching coefficients of
$\bar\mu^2$:
\[
    \dot{C}_t = -2\kappa B_t + \frac{2A_tC_t}{\lambda} + \frac{B_t^2}{\lambda} +
    \frac{2B_tC_t}{\lambda}, \qquad C_T = 0.
\]
These replace \eqref{eq:sr-riccati-quadratic} (and, with $Q,P$ in place of
the scalar $c_1/2,c_2/2$ and matrix products throughout, the analogous
system in Chapter~8's \eqref{eq:alpha-riccati-system} should be re-derived
under the same correction; this is flagged as remaining work in
Chapter~10, \S10.4, since the matrix case was not re-derived in this
revision).

\subsection{Linear Terms}

Matching coefficients of $x$:
\[
    \dot{D}_t = \kappa D_t + \frac{A_tD_t}{\lambda} + \frac{B_tD_t}{\lambda} -
    (A_t+B_t)b_t, \qquad D_T = 0,
\]
which adds the $B_tD_t/\lambda - B_tb_t$ master-equation contribution to the
previously stated equation. Matching coefficients of $\bar\mu$:
\[
    \dot{E}_t = -\kappa D_t + \frac{A_tE_t}{\lambda} + \frac{B_tD_t}{\lambda} +
    \frac{B_tE_t}{\lambda} + \frac{C_tD_t}{\lambda} - (B_t+C_t)b_t,
    \qquad E_T = 0.
\]
This confirms the open item flagged in an earlier revision: $\dot{E}_t$ does
indeed depend on $C_t$, through the $C_tD_t/\lambda$ term. Matching the
constant term, which collects the diffusion contributions
$\partial^2_{xx}V=-A_t$, $\partial^2_{x\bar\mu}V=-B_t$,
$\partial^2_{\bar\mu\bar\mu}V=-C_t$ together with the remaining
$\bar{b}\cdot\partial_{\bar\mu}V$ pieces:
\[
    \dot{F}_t = \frac{D_t^2}{2\lambda} + \frac{D_tE_t}{\lambda} -
    (D_t+E_t)b_t - \frac{\sigma^2+\sigma_0^2}{2}A_t - \sigma_0^2B_t -
    \frac{\sigma_0^2}{2}C_t, \qquad F_T = 0.
\]
The full system $(A_t,B_t,C_t,D_t,E_t,F_t)$ was verified to make the PDE
residual vanish identically under direct symbolic substitution, and was
integrated numerically (Section~\ref{sec:sr-numerical}) with no blow-up
across the full original parameter range, including the originally intended
$c_2=5.0$. As in the previous revision, when $b_t\equiv0$ (Chapter~8's
alpha-generation model has no exogenous bailout term), $D_t=E_t\equiv0$, but
$C_t$ and $F_t$ no longer vanish in general -- the simplification used in
Chapter~8 should therefore be re-checked against this corrected system as
part of the matrix-case re-derivation noted above.

\subsection{Proof that $v_t$ is Independent of $b_t$}

From the variance dynamics \eqref{eq:sr-variance-dynamics}, the variance $v_t$
satisfies a linear ODE whose coefficients depend only on $A_t, \kappa, \lambda,
\sigma$, and $\sigma_0$. None of these depend on the bailout policy $b_t$.
Therefore, the government's choice of $b_t$ affects the mean $\bar\mu_t$ but not
the variance $v_t$. This justifies the separation in the government's objective
used to derive \eqref{eq:sr-optimal-bailout}.

\subsection{Government's Riccati Equation (Chapter 7)}

The government's value function takes the form $V(t,\bar\mu) = -\frac{1}{2}\Pi_t
\bar\mu^2 - \delta_t$. Substituting into the HJB equation for the government's
control problem (with dynamics \eqref{eq:sr-mean-dynamics} and cost
\eqref{eq:sr-gov-objective}) and matching coefficients of $\bar\mu^2$ yields
\eqref{eq:sr-gov-riccati}.

\subsection{Explicit Bound for the Propagation of Chaos Constant (Chapter 7)}

The bound for $C$ in Proposition~\ref{prop:sr-propagation-of-chaos} follows from
the synchronous coupling estimate of Chapter 3. Specifically, let
$\Delta X_t^i = X_t^{i,N} - X_t^{i,*}$ be the difference between the $N$-bank
state and the limiting McKean--Vlasov state under the same noises. Then, by the
same It\^o-expansion argument as in \Cref{sec:appA-contraction},
\[
    \mathbb{E}\big[|\Delta X_t^i|^2\big] \leq \int_0^t e^{2L_b(t-s)} L_b\,
    \mathbb{E}\big[W_2(\mu_s^N, \mu_s^*)^2\big]\, ds.
\]
Summing over $i$ and using $\mathbb{E}[W_2(\mu_t^N,\mu_t^*)^2] \leq
\frac{1}{N}\sum_{i=1}^N \mathbb{E}[|\Delta X_t^i|^2]$ yields the claimed bound via
Gronwall's inequality, with the explicit constant
\[
    C \leq \frac{4}{\kappa_{\min}}\left( \sigma^2 + \sigma_0^2 + \frac{1}{\lambda^2}
    \sup_t (A_t+B_t)^2 \right) e^{2L_b T},
\]
where $\kappa_{\min} = \inf_t(\kappa + A_t/\lambda) > 0$ and $L_b = \kappa +
\sup_t(A_t + |B_t|)/\lambda$, as stated in Chapter 7.
